\documentclass[11pt,a4paper,sans]{moderncv}
\usepackage{fontawesome5}
% The file must come before \moderncvstyle because
% default-socialicon-color is used to set the color of social icons

%-------------------------------------------------------------------------------
%                color scheme definition
%-------------------------------------------------------------------------------
\definecolor{primary}{HTML}{1447e6}
\colorlet{color1}{primary}
\colorlet{default-socialicon-color}{color1}

%-------------------------------------------------------------------------------
%                end of color scheme definition
%-------------------------------------------------------------------------------

\moderncvstyle{casual}
\moderncvcolor{blue}
\usepackage[french]{babel}
\usepackage[top=1cm,bottom=2cm,left=1cm,right=1cm]{geometry}
\usepackage{csquotes}
\usepackage[T1]{fontenc}
\usepackage[utf8]{inputenc}
\usepackage{textcomp}
\newcommand{\lien}[3][blue]{\href{#2}{\color{#1}{#3}}}
\nopagenumbers{}

%-------------------------------------------------------------------------------
%                color scheme definition
%-------------------------------------------------------------------------------
\colorlet{lastnamecolor}{primary}
\colorlet{namecolor}{lastnamecolor}
\colorlet{firstnamecolor}{lastnamecolor}
\colorlet{headrulecolor}{color1}
\colorlet{titlecolor}{color1}
\colorlet{addresscolor}{primary}
\colorlet{quotecolor}{color1}
\colorlet{pictureframecolor}{color1}
%-------------------------------------------------------------------------------
%                end of color scheme definition
%-------------------------------------------------------------------------------


%%%%%%%%%%%%% Beginning of \makelettertittle modification %%%%%%%%%%%%%%%%%%

\makeatletter
%-------------------------------------------------------------------------------
%                letter style definition
%-------------------------------------------------------------------------------
% commands
\renewcommand*{\recomputeletterlengths}{%
    \recomputecvlengths%
    \setlength{\parskip}{6\p@}}

\providecommand*{\subject}[1]{\def\@subject{#1}}

\renewcommand*{\makelettertitle}{%
% recompute lengths (in case we are switching from letter to resume, or vice versa)
    \recomputeletterlengths%
% sender contact info
%\hfill%
    \begin{minipage}{.5\textwidth}%
        \raggedright%
        \addressfont\textcolor{primary}{%
                {\bfseries\upshape\@firstname~\@familyname}\@firstdetailselementfalse%
            \ifthenelse{\isundefined{\@addressstreet}}{}{\makenewline\addresssymbol\@addressstreet%
            \ifthenelse{\equal{\@addresscity}{}}{}{\makenewline\@addresscity}}%
            \collectionloop{phones}{\makenewline\csname\collectionloopkey phonesymbol\endcsname\collectionloopitem}%
            \ifthenelse{\isundefined{\@email}}{}{\makenewline\emailsymbol\emaillink{\@email}}%
            \ifthenelse{\isundefined{\@homepage}}{}{\makenewline\homepagesymbol\httplink{\@homepage}}%
            \ifthenelse{\isundefined{\@extrainfo}}{}{\makenewline\@extrainfo}}%
    \end{minipage}\vspace{2em}% << Add space with \vspace, have actual paragraph break

% recipient block        V--V Adjust to fit text and align as desired
    \hfill\begin{minipage}[t]{.5\textwidth}       % << \hfill now works
              \raggedright%
              \addressfont%
              {\bfseries\upshape\@recipientname}\\%
              \@recipientaddress\\[2em]%
              \@date\\[1.5em]%
    \end{minipage} % << Another paragraph break here needed

% date
    \hfill  % << This wasn't doing anything good before, probably not now either.
    \null\\[2em]%
%\raggedright % << Remove ragged right from definition
    \ifthenelse{\isundefined{\@subject}}{}{%
        \begin{tabular}{ p{0.08\textwidth} p{0.92\textwidth} }
            Objet: & \raggedright \@subject
        \end{tabular}
        \\[1.5em]}%
% opening
%\raggedright % << Remove ragged right from definition
    \@opening\\[1.5em]%
    \hspace{0pt}\par\vspace{-\baselineskip}\vspace{-\parskip}}

\renewcommand*{\makeletterclosing}{
    \@closing\\[3em]%
    {\bfseries \@firstname~\@familyname}%
    \ifthenelse{\isundefined{\@enclosure}}{}{%
        \\%
        \vfill%
        {\color{color1}\itshape\enclname: \@enclosure}}}
\makeatother%
%%%%%%%%%%%%%%%%% End of \makelettertittle modification%%%%%%%%%%%%%%%%%%%%%%%%%%%

\patchcmd{\makeletterclosing}             %%% Flushing right the signature
{\@closing}
{\raggedleft \@closing}
{}{}
\makeatother

% Informations personelles
\name{Amadou}{Sall}
\address{4 rue Claude Berri}{69800 Saint-Priest}{} % pas nécessaire de mettre code postale et ville
\phone[mobile]{+33~6~95~60~68~64}        % l'option peut être "mobile" (default), "fixed" or "fax"
%\email{ahasall@yahoo.fr}
%\social[github]{ahasall} % option = "twitter", "github" ou linkedin
%\social[linkedin]{ahasall}
%\social[twitter]{ahasall}


% Destinataires et logement concerné.
\newcommand{\nomLocataire}{Monsieur Quentin Bertrand\\Madame Chloé Piger}
\newcommand{\adresseDestinataire}{286 avenue de Lardenne\\Appartement 13\\31100 Toulouse}
\newcommand{\adresseLogement}{286 avenue de Lardenne, appartement 13, 31100 Toulouse}

\title{Régularisation des charges locatives du 9 novembre 2023 au 30 juin 2025}
\recipient{\nomLocataire}{\adresseDestinataire}
\date{Saint-Priest, le 27 octobre 2026}
\subject{Régularisation des charges locatives des exercices 2023-2024 et 2024-2025}
\opening{Madame, Monsieur,}
\closing{Cordialement,}

\begin{document}
    \hypersetup{
        pdftitle={Régularisation des charges locatives du 9 novembre 2023 au 30 juin 2025},
        pdfauthor={Amadou Sall},
        pdfsubject={Régularisation des charges du logement situé 286 avenue de Lardenne, appartement 13, 31100 Toulouse},
        pdfkeywords={charges, régularisation, 2023, 2024, 2025, Toulouse}
    }
    \makelettertitle

    Conformément à l'article 23 de la loi n\textsuperscript{o}~89-462 du 6 juillet 1989, je vous adresse le décompte de régularisation des charges locatives du logement situé \adresseLogement{} et des deux parkings couverts (lots 27 et 28), pour la période du \textbf{9 novembre 2023 au 30 juin 2025}.

    Les régularisations suivent les exercices comptables de la copropriété, du 1\textsuperscript{er} juillet au 30 juin. Les charges de copropriété correspondent à la seule part récupérable indiquée par le syndic (décret n\textsuperscript{o}~87-713 du 26 août 1987). Les charges de copropriété et la taxe d'enlèvement des ordures ménagères (TEOM) sont calculées au prorata des jours d'occupation.

    \textbf{Exercice du 9 novembre 2023 au 30 juin 2024} (235 jours d'occupation sur 366)
    \begin{center}
        \small
        \begin{tabular}{p{0.42\textwidth} p{0.28\textwidth} r}
            \hline
            Charges de copropriété récupérables & 828,59 $\times$ 235 / 366 & 532,02~\texteuro \\
            TEOM 2023, du 9 novembre au 31 décembre & 213,00 $\times$ 53 / 365 & 30,93~\texteuro \\
            TEOM 2024, du 1\textsuperscript{er} janvier au 30 juin & 221,00 $\times$ 182 / 366 & 109,90~\texteuro \\
            \textbf{Total des charges réelles} & & \textbf{672,85~\texteuro} \\
            Provisions versées & 51,33 + 7 $\times$ 70,00 & $-$~541,33~\texteuro \\
            \hline
            \textbf{Solde dû} & & \textbf{131,52~\texteuro} \\
            \hline
        \end{tabular}
    \end{center}

    \textbf{Exercice du 1\textsuperscript{er} juillet 2024 au 30 juin 2025}
    \begin{center}
        \small
        \begin{tabular}{p{0.42\textwidth} p{0.28\textwidth} r}
            \hline
            Charges de copropriété récupérables & exercice complet & 810,50~\texteuro \\
            TEOM 2024, du 1\textsuperscript{er} juillet au 31 décembre & 221,00 $\times$ 184 / 366 & 111,10~\texteuro \\
            TEOM 2025, du 1\textsuperscript{er} janvier au 30 juin & 225,00 $\times$ 181 / 365 & 111,58~\texteuro \\
            \textbf{Total des charges réelles} & & \textbf{1\,033,18~\texteuro} \\
            Provisions versées & 12 $\times$ 70,00 & $-$~840,00~\texteuro \\
            \hline
            \textbf{Solde dû} & & \textbf{193,18~\texteuro} \\
            \hline
        \end{tabular}
    \end{center}

    Le montant total restant dû au titre de ces deux exercices s'élève donc à \textbf{324,70~\texteuro{}} (131,52~\texteuro{} + 193,18~\texteuro{}).

    Les pièces justificatives (décomptes du syndic détaillant les charges et leur mode de répartition entre les lots, avis de taxe foncière) restent à votre disposition pendant six mois à compter de cet envoi.

    Je reste à votre disposition pour toute question.

    \makeletterclosing
\end{document}
